\chapter{todo???Is Surrogacy a Special Kind of Labour That Should Not Be Treated According to Market Norms?}
\chaptermark{Is Surrogacy a Special Kind of Labour}
\chapterauthor{todo,
\textit{todo}}

\renewcommand*{\thesection}{\arabic{section}.}
\renewcommand*{\thesubsection}{\arabic{section}.\arabic{subsection}:}

\begin{quote}
    This paper examines whether
surrogacy constitutes a special kind of labour that means it should not
be treated according to market norms/be commercialised. I evaluate
arguments from Elizabeth Anderson who I believe fails to explain why
surrogacy should not be treated according to market norms due to her
claim that there is an essential feature of gestation that renders it
unable to be commercialised; I draw on Cecile Fabre's claim that the
demands of surrogacy are comparable to those found in other professions
we deem to be permissible. Yet, I argue that Fabre's claim does not
render commercial surrogacy permissible; turning to Debra Satz's account
of the asymmetry thesis, I argue that while there is nothing
intrinsically special about surrogacy, it should still be treated
asymmetrically because of the social context in which it occurs. I argue
in particular that commercialising surrogacy risks reinforcing
preexisting patriarchal norms and the idea that women's bodies should be
controlled, which I believe contributes to gender inequality. I consider
some critiques against the success of Satz's asymmetry thesis,
particularly focussing on concerns about paternalism and
overgeneralisation, yet I defend a revised version of the asymmetry that
emphasises the broader social harms of commercial surrogacy to women
uninvolved in the sector at all. I ultimately conclude that while
commercial surrogacy is not inherently objectionable, it should still
not be treated according to market norms due to its harmful role in
sustaining gender inequality in our society. 
\end{quote}

\section*{Introduction}

Surrogacy is
the process in which a woman carries a child to term for another person
or couple. In commercial surrogacy, this arrangement involves financial
compensation to the surrogate and is governed by an enforceable, legal
agreement; this differs from altruistic surrogacy where the surrogate
does not receive monetary compensation and often is motivated by
personal considerations. Commercial surrogacy remains a highly
controversial practice, being illegal in the UK, Canada, Australia, and
many countries across Europe and Asia. Discussion on commercial
surrogacy typically centers on the potential exploitation of surrogates
(particularly those in low-income regions), the commodification of
children, and the surrogates' own control over their medical and
personal decisions. Some feminist theorists hold that there is something
inherent about the act of bringing a child to term which renders
surrogacy a special kind of labour that should not be commercialised. I
am particularly interested in identifying this essential property which
excludes surrogacy from being able to be treated according to market
norms, that is, being treated according to the rules and expectations
governing commercial exchanges. 

The view that there is something special
about surrogacy that means it should not be commercialised is what I
will be referring to as the asymmetry thesis, which in other words, is
the argument that surrogacy should be treated asymmetrically to other
forms of labour; I will be considering Elizabeth Anderson's asymmetry
thesis in particular. In this paper, I will argue that Anderson fails to
show that surrogacy constitutes a special kind of labour, yet I will
ultimately endorse a version of the asymmetry thesis motivated by
different considerations than this view. I will begin by presenting
Anderson's claim that commercial surrogacy is a special kind of labour
and should be treated distinctly to other forms of labour. 

I argue
that her argument comes under fire from Fabre's rejection that surrogacy
should be treated asymmetrically, and I maintain that there is nothing
essentially unique to surrogacy as a line of work. I provide some
objections to Fabre, yet will show that while surrogacy is not
inherently objectionable, there is still something intuitively
disagreeable about selling reproductive labour. I present Satz's version
of the asymmetry thesis as an explanation for the intuition that
surrogacy should not be commercialised. In evaluating Satz's argument, I
point out that her claim risks paternalism by implying that surrogates
cannot freely choose their line of work in patriarchal society. I also
show that Satz seems to have missed the fact that many stereotypically
gendered lines of work exist and are able to be treated according to
market norms. Despite these critiques, I argue that Satz's asymmetry
thesis still succeeds, and I show that these counters do not undermine
her argument. I ultimately conclude that what makes commercial surrogacy
objectionable is not it being `special' in nature, but rather its
contributions to essentialist norms concerning the stereotypical social
roles of women in a patriarchal society. Surrogacy should therefore not
be treated according to market norms. 

\addtocounter{section}{1}
\subsection{Anderson's Case Against
Commercial Surrogacy}

Commercial surrogacy involves a prospective
parent/parents paying a woman to carry a child, compensating her for
both the labour of pregnancy and for waiving her parental rights towards
the child. Anderson (1990) argues that commercial surrogacy is
objectionable, as it is trying to transform what is specifically a
labour done by women into a commodity which follows economic norms of
typical labour markets. She contends that there is an asymmetry between
other types of labour and surrogacy: carrying a child to term for
another is a service that would be wrong to sell or buy, and contract
pregnancy wrongfully allows for the extension of the market into the
``private'' sphere of reproduction. 

Her argument for this claim begins
by highlighting respect and consideration as two modes of valuation
which are disregarded by commercial surrogacy. She defines respect for
someone as treating them as an autonomous person with rational
interests. While to treat someone with consideration is to respond to
her and her emotional relations with others with sensitivity, avoiding
any manipulation or belittling of these feelings, especially for one's
own benefit. Applying market norms to reproductive labour violates both
of these in three ways: (Anderson 1990, 81): 

\vspace{-1em}
\textbf{\begin{enumerate}
\renewcommand{\labelenumi}{(\roman{enumi})}
\item
  Surrogates are alienated
as they are required to suppress their emotional connection with the
child. 
\item 
    Surrogates are degraded as the legitimacy of their personal
evolving perspective on their own pregnancy is denied, or is demoted to
an amount of money. 
\item 
    Surrogates are exploited as their emotional
needs are not considered as relevant, yet may also be manipulated to
encourage them to undergo surrogacy. 
\end{enumerate}}

\noindent These three ways of denying
respect and consideration, Anderson argues, subsequently treat women's
bodies as a commodity. The key idea that surrogacy firms miss when
commercialising surrogacy is the role of emotions which are typical of
non-commercialised pregnancies. Pregnancy is not merely a biological
process---it is a social practice where parents prepare themselves to
welcome new life (Anderson 1990, 81). Without regard for this crucial emotional
aspect, commercial surrogacy entails `hiring' a woman's womb in exchange
for payment. In turn, the surrogate is reduced to an object for mere use
rather than a person worthy of respect. 

For Anderson, commercial
surrogacy disregards natural emotional processes and therefore alienates
surrogates. Anderson holds that it is common that many women develop
feelings of love towards the child they carry, yet that there is no
place for this reaction in commercial surrogacy, since the surrogate
mother agrees not to form a parent-child connection as her contract
states she must give up the child to the clients. The potential risk of
this being unbearably painful is so common that 10\% of women suffer
grief so intense that they require professional help (Anderson 1990,
83). In turn, the surrogate is deliberately required to alienate herself
from her love for the child, which can be an incredibly demanding task.
Failure to give up the child would be treated as a breach of contract or
an incomplete transaction. Yet the certainty promised to prospective
parents---that they can expect to receive a child---does not exist for
the surrogate. It is not certain that the surrogate will be able to bear
the psychological effects that may come with relinquishing the child.
She is expected to take on the risks of transforming her emotional
perspective on herself with the same detachment as a trader brings to
fluctuations in a product's price (Anderson 1990, 84). This is degrading, since
an instinctive human emotion is not taken into account, and so in some
way also denies the humanity of the surrogate, reducing her to a
vessel whose sole function is to produce children. The emotional ties
which she may establish with the child are ignored. 

In light of
Anderson's claim, it is also important to recognise that the emotional
aspect of surrogacy is extraordinarily difficult to take into account
for compensation. Pregnancy, while extremely demanding, can also be
incredibly varied. Consider two surrogates: one experiences no major
health complications and feels little emotional distress when giving up
the child. The other however, develops high blood pressure, gestational
diabetes, infections, and severe postnatal depression resulting from
emotional turmoil. Both these surrogates will have all their distinct
medical expenses covered and receive support for their physical needs.
Yet, a clear difference remains. The second surrogate has dealt with a
much harder emotional journey. Despite this, the base price of her
labour, excluding additional medical expenses, will be similar to that
of the first. Consequently, she has not been sufficiently compensated.
Anderson points out that the prospective parents are required to pay for
all medical expenses, with the exception of long-lasting mental health
impacts on the surrogate mother; therapy post pregnancy is not accounted
for in the cost of surrogacy. (Anderson 1990, 83). I also want to
emphasise that emotional labour is not a quantifiably measurable
commodity, which makes fair compensation challenging---it is much harder
to put a price on a surrogate's unique emotional experience in
comparison to her hospital bill. Yet, a uniform cash sum completely
ignores the deep emotional implications of pregnancy and its nature as
extremely personal. I can anticipate the response that surrogates are
paid generously for the very fact that surrogacy is an emotionally
taxing role, a standardised paycheque makes assumptions and generalises
emotional experiences, which will often vary greatly. In turn, Anderson
argues that this is exploitative, as it fails to correctly take the
surrogate's labour and personal needs into account. 

\subsection{Objections from Fabre}

Fabre (2006) argues that Anderson is wrong to conclude that
surrogacy is a different type of labour, as the demands and dangers that
make it `special' can be likewise attributed to other professions. There
are multiple physically demanding non-typical jobs, some of which could
be considered even more risky than surrogacy. For instance, astronauts
are subject to suit malfunctions, loss of bone density and muscle
atrophy. Surrogacy likewise deals with unexpected health risks, such as
postpartum haemorrhages, long-term fertility effects and placental
complications. If people, despite these implications, can choose the
incredibly difficult job of being an astronaut, we should agree that
surrogacy, which is also considered to be an unconventional line of
work, should be permissible and treated similarly (Fabre 2006, 194--195).

Anderson may urge that the emotional demands rather than the physical
are what sets surrogacy apart, yet I anticipate that Fabre could also
reject this. Social work for instance, also has the ability to deeply
alter one's mental state. Social workers are often exposed to extremely
emotionally taxing situations due to dealing with disturbing cases.
Their lives could certainly be altered negatively and permanently by
such trauma. Yet, I take that Fabre would argue that we should not ban
social work, even if it is a profoundly difficult job, as we assume that
social workers fully recognise and consent to being in such situations.
It seems strange to deny surrogates the same autonomy to consent to a
job with similar emotional risks. Surrogates are aware that they are
signing contracts that consent to the possibly painful experience of
relinquishing a child. Rejecting this awareness seems paternalistic, as
we should assume they have autonomy to willingly choose to be
surrogates (Fabre 2006, 212). We do not think twice about men who willingly
sign up to the army for example, where one has a high chance of dying in
battle; we confidently assume they have capacity to make that choice
themselves, and we should do the same for surrogates.

Regarding fair
compensation, it may be the case that some surrogacy firms do not
account for emotional complications and thus do not fairly pay
surrogates. Yet, Fabre could respond that this shortcoming does not
render surrogacy as something that cannot be commercialised. The social
worker and the astronaut alike would also be unfairly paid if their
emotional complications were not compensated for. These high risk jobs
do however take complications into account; astronauts are given hazard
pay for instance, which is an additional income for working in dangerous
conditions. Surrogacy also has the potential to be satisfactorily
compensated for, and perhaps surrogacy firms are merely not doing enough
currently. This does not imply there is something essential which
separates surrogacy from other jobs, but that there is something unfair
within the models that currently exist, which can certainly be amended.

\subsection{Responses to Fabre}

I can imagine a response from Anderson in
defence of her claim possibly taking this form: the emotional aspect of
surrogacy is uniquely significant because it involves the dismissal of a
personal and intimate biological process. I imagine Anderson would argue
that the emotional aspect is distinct not due to being merely a
difficult job, yet because it involves the intimate and intense
experience of creating life within one's body. Relinquishing a child one has carried to term can lead to extreme emotional challenges that
greatly differ from challenges in other roles, such as social work.
Social workers may encounter anguish, yet it is not their child that
they feel this anguish about. The emotional distress of surrogates is
magnified by the fact that they must give up their own parental
instincts; regardless of expectations to repress this, giving up a child
one has formed a connection to over nine months is uniquely painful and
unnatural according to Anderson (1990, 81--82). 

However, I argue Anderson
is wrong to assume that this life-altering, unbearable emotional
distress is an inherent or universal experience for surrogates. Not
every surrogate has a maternal bond or deep attachment to the child she
carries; many are able to relinquish the child without much difficulty.
In fact, most surrogates enter surrogacy contracts with altruistic
motives, and take great pride and satisfaction in doing this service. I
also believe that placing so much emphasis on a natural maternal
instinct seems unsavoury and almost stereotypical of women itself. Many
women struggle to feel maternal instincts and desire for children; some
even abort them. Arguing that a deep emotional connection is what
pregnancy naturally or normally involves sets an expectation for women
to be naturally nurturing or caregiving; this stigmatises women who lack
these emotions, deeming them as selfish and abnormal. The reliance on a
`bond' that sets surrogacy apart is unsuccessful and reinforces
traditional stereotypes for women. The existing diversity of women's
experiences with pregnancy therefore weakens Anderson's claims that
surrogacy is special, as surrogacy does not necessarily lead to severe
emotional harm. 

While I do agree that there is nothing essentially wrong
with surrogacy, I also maintain that there is something that makes us
feel uneasy about commercialising women's bodily functions, a point
which Fabre's account misses. While there may not be anything
essentially wrong with surrogacy, we may still share the sentiment that
there is something that indeed does set surrogacy apart. 

\renewcommand*{\thesection}{\arabic{section}:}
\section{Debra Satz's Asymmetry Thesis}
  
I argue that this intuition can be answered by Satz
(1992), who defends the asymmetry thesis on different grounds. In this
section, I will provide an account of Satz's argument before my analysis
of its success in sections 3.1 and 3.2, in which I deem this version of
the asymmetry thesis a compelling explanation for why surrogacy should
not be commercialised. 

Satz rejects Anderson's claims that surrogacy is
a special type of labour, yet recognises that we should treat surrogacy
differently as it exists in our patriarchal society, and thus reinforces
a traditional gender-imbalanced division of labour. She claims that
there should be an asymmetry between the treatment of surrogacy and
other labour markets, as a patriarchal climate renders surrogacy markets
to be much more problematic, since gender inequality already exists in
labour markets. Even if a woman chooses to work as a surrogate, we must
recognise that this choice exists with unequal opportunities for women
in the background (Satz 1992, 124). My focus will be on two of Satz's
objections to contract pregnancy: firstly, contract pregnancy allows
increased access and control over women's bodies and secondly, that it
risks reinforcing pre-existing patriarchal norms women are subject to.

Satz asserts that there is nothing essentially, or morally wrong with
reproductive labour, yet it should still be treated distinctly, as
contract surrogacy upholds gender inequality by putting women in a
position to be controlled through increased access to their sexuality
and body (Satz 1992, 124). She points to the difference between surrogacy and
other forms of reproductive labour, such as men donating sperm. In the
latter, sperm donation does not give anyone any control over the
sexualities and bodies of the men; once a man donates his sperm, his
bodily involvement in the reproductive process comes to an end. There is
no ongoing physical or behavioural control or restriction to him. Unlike
this, surrogacy entails selling control over one's body (Satz 1992, 125). To
give birth to a healthy child, women must be restricted from consuming
certain foods and drinks, and avoiding strenuous activities, and medical
decisions may also be made on her behalf. This restriction lasts for a
prolonged period of time, as the surrogate's body remains under
continuous monitoring and restriction for months. In essence, this
creates a situation where other people (the parents, doctors, or
contractors) have some sort of control over the woman's body. There are
other likewise controlling jobs such as athletes, whose bodies are also
restricted to having specific diets and exercise routines (possibly for
years), yet Satz stresses that the restriction on surrogates is still
underlined by something different. 

This difference is the type of
control that surrogacy agencies have over surrogates. It is not
inherently objectionable that someone may have control over someone
else's body through a contract, yet it is objectionable that someone has
control over a woman's body because women's bodies have been controlled
historically (Satz 1992, 125). Underlying gender inequality is not similarly
tied to the work of athletes. While they may undergo strict diets or
exercise routines that we may consider to be invasive and difficult,
their gender has nothing to do with this restriction. Their profession
does not have to do with their sex---anyone can be an athlete. Not
everyone can, however, be a surrogate. The job relies entirely on a
woman's sex and sexual capacities. The surrogate's reproductive journey
is controlled in some way by her agency, and this power dynamic seems to
reinforce the misogynistic, continuous control over women's bodies.

Secondly, commercial surrogacy, when set in the context of our current
society, also reinstates the gender stereotype that women exist solely
for the purpose of child rearing (Satz 1992, 127). The purpose of working as
a surrogate is to have children as a service; in a patriarchal world,
surrogacy therefore reinforces the idea that women are merely wombs or
vessels for creating children, as the identities of surrogates are
reduced to their reproductive capacities. When commercialising surrogacy
in a society such as ours that has these misogynistic qualities, the
women who work as surrogates are degraded as `baby machines', whose
value lies in their reproductive abilities. It is also the case that the
reinforcement of such stereotypes may lend itself to the views of the
surrogates, where even they view themselves as having value solely due
to how well they perform their profession as child rearers. Satz argues
that this in turn is harmful to the feminist project, as patriarchal
society does not view these women as autonomous, multifaceted
human-beings and rather as objects with `wombs for hire'. This pushes
the reductive stereotype that the primary function of a woman is tied to
reproduction and thus perpetuates the social subordination of women
(Satz 1992, 127). 

\addtocounter{section}{1}
\subsection{Objections to Satz}

I can imagine one pointing out that
Satz's argument may be deemed as paternalistic. While it tries to
protect women from gender inequality, it seems to suggest that women
cannot make any choices for themselves and are completely helpless at
the hands of patriarchal society. It does not feel aligned to the
feminist position in the first place to deny that women have the ability
to consciously choose their line of work and fully know what their job
entails. I can imagine that many surrogates take pride in what they do,
arguing that their job is fulfilling because they can help so many
individuals who are unable to carry their own biological children, such
as women with ovarian cancer, and gay couples. Denying this is telling
fully functional, adult women that their choices are mistaken and that
their consent is invalid; it tells women that they simply `do not know
any better' as victims of society's control. Even if we accept that
perhaps for the most part, those who choose to be surrogates do suffer
from gender inequality, banning commercial surrogacy for the `greater
good' still risks denying at least some women the opportunity to do work
they value. 

I also want to highlight multiple other labours that could
wrongly imply things about women. For instance, domestic labour has been
historically performed by women, and women are often stereotyped as
being suited for such labour on the assumption that they are inherently
nurturing (2025). Taking Satz's line of argument, I hold that we could
make a similar point about how for example, women working as cleaners
can reinforce the misogynistic idea that women are made for housework. I
further want to extend one of the previous arguments Satz makes in
emphasising the even greater exploitation and inequity between employers
and women who are socioeconomically disadvantaged, pointing to how many
women who work as caretakers do come from such underprivileged
backgrounds. If we follow Satz's asymmetry thesis, perhaps other
stereotypically feminine jobs threaten the feminist project in the same
way surrogacy does. Yet, it seems absurd that we should consider
rendering these jobs impermissible in the same way commercial surrogacy
is argued to be objectionable. 

\subsection{A Defence of Satz}
  
I argue that
the asymmetry thesis can still be defended despite these two critiques.
While it might seem paternalistic due to its apparent dismissal of the
capacity for consent surrogates as adult women have, I argue that it is
not. Much like the liberal feminist idea that women should be free to
make any choice without moral consideration, I take the defence of
surrogacy here to be problematic. 

I believe that the same arguments used
to reject defences of pornography could meaningfully apply here. Many
feminist theorists condemn pornography not only because of the
exploitation of the actors themselves, but because of the harmful
secondary effects it has on women who are uninvolved in its production.
According to Dworkin and MacKinnon, viewers of pornography are
influenced to have harmful and violent attitudes towards women in their
actual lives because of the depictions of extreme violence and female
subordination in pornography (1988, 25). Moreover, as an industry made
`for men by men', women are often depicted to be in compromising,
subordinate positions, and shown to `serve' the interests of their male
partners. This therefore upholds the patriarchal idea that women are
sexual objects to be used to achieve the ends of male pleasure, which is
something that affects all women (1989, 206--209) and thus renders
pornography impermissible. I argue that this type of argument should be
applied in response to the accusations of paternalism against Satz's
asymmetry thesis. Surrogacy can also be seen as something that is
harmful not only for the surrogates involved, whose reproductive labour
directly reduces them to vessels for their clients' children, yet
reinforces this stereotype in general, which affects the societal value
of women as a whole. Even if we assume that some women do find
fulfilment in being surrogates, the very existence of this line of work
in a patriarchal climate is detrimental to all women, extending to those
who are not involved in commercial surrogacy at all. I argue that
therefore it is justified to reject commercial surrogacy, as there is
still great harm being done to women even if we grant that some women
can truly enjoy being surrogates. I do not believe it is paternalistic
to admit this, and I believe that the consideration of secondary effects
onto all women (which I take to be missed out by even Satz who focuses
on the subordination of the surrogates themselves) strengthens this
version of the asymmetry thesis. 

Regarding the second critique, I argue
that comparisons of surrogacy to other typically feminine jobs are
disanalogous. While surrogacy and other stereotypically feminine jobs
may both wrongly imply things about the women who perform them,
caretaking professions such as cleaning are not defined strictly on the
basis that they are done by biological women. Working as a cleaner may
be surrounded by stereotypes that society has put upon the women in this
profession, yet being a woman does not define or establish what this job
entails. Anyone can work as a cleaner regardless of the fact that the
job is stereotypically feminine; being a certain sex is not tied to what
the profession entails. We can also see this in stereotypically male
jobs, such as being in the military, where these jobs similarly do not
entail or require the worker to be a biological man. Alternatively,
surrogacy is defined by and tied to being female entirely, as child-rearing is exclusive to biological women. The male sex cannot
participate in this line of work at all, and so it is a requirement that
the workers are biological women, able to reproduce. While we could
have the capacity to eliminate misogyny towards stereotypically female
lines of work, as stereotypes attributed to them are solely societal, I
maintain that we cannot say the same for surrogacy. Since surrogacy is
selling child-rearing and is exclusively performed by women, it is much
harder to escape female objectification and exploitation in a
patriarchal world. Women have historically been seen as `baby-making'
sexual objects which reduces the value of women to their reproductive
capabilities. This type of misogyny is much harder to shake, and so the
asymmetry thesis holds. 

Regarding the status of women's reproductive
labour, I believe that the asymmetry thesis is a convincing argument to
explain what makes surrogacy something that cannot and should not be
commercialised. I recognise that there are also other negative effects
of commercial surrogacy that do not have to do with its existence within
the unequal social relations of gender which I have not considered in
this paper (such as the rearing and selling of children), but I hold
that the upholding of gender inequality is nevertheless a principal
reason to reject permitting commercial surrogacy. 

\section*{Conclusion}

Surrogacy
must still be treated differently to other forms of labour, even if
there is nothing essentially morally wrong with it. While I have argued
that Fabre succeeds in undermining Anderson's defence of the asymmetry
thesis on grounds that surrogacy is fundamentally problematic, I hold
that we should still treat surrogacy differently. I have shown how
Satz's asymmetry thesis successfully explains this on the grounds of the
unequal social context in which commercial surrogacy occurs, which in
turn reinforces gendered stereotypes which reduce women to their
reproductive abilities and stereotyped roles. I have responded to
criticisms of paternalism I have anticipated against Satz, as well as
showing that while there are other traditionally feminine professions,
surrogacy is distinct. Surrogacy is defined by the biological capacities
of biological women, and therefore cannot be considered in isolation
from misogynistic structures which often appeal to those same biological
traits in justifying the oppression of women. Therefore, while surrogacy
is not a special labour, it should not be treated according to market
norms, as doing so reinforces the gender inequality persistent in our
society. 

\refsection

\begin{hangparas}{\hangingindent}{1}
    Anderson, E. S. (1990). Is Women's Labor a Commodity?. Philosophy \& Public Affairs, 71--92.13. 

    Fabre, C. (2006). Surrogacy
Contracts. In Whose Body is it Anyway? Justice and the Integrity of the
Person. Oxford University Press. 

    MacKinnon, C. A, and Dworkin, A.
(1988). Pornography and Civil Rights: A New Day for Women's Equality.
Minneapolis, MN: Organizing Against Pornography. 

    MacKinnon, C. A.
(1989). Towards a Feminist Theory of the State. Harvard University
Press. 

Satz, D. (1992). Markets in women's reproductive labor.
Philosophy \& Public Affairs, 107--131.14. 

    UN Women. (2025). FAQs: What
is unpaid care work and how does it power the economy?. UN Women.
\end{hangparas}